% !TeX root = ../main.tex

\section{Introduction}


The copyright should be honored, both in literature, required by governments. and in code, evaluated in many legal cases. 

There is emerging conflicts unresolved. Model providers reluctant to provide. transparency issue exists. 


MIA is a long-studied  topic. Initially, MIA was extensively studied in machine learning. 

With the rise of LLMs, MIA was also introduced. 

Yang 


Tirumala et al. analyze the memorization dynamics of different parts of speech and find that models memorize nouns and numbers first, that larger models can memorize a larger portion of the data before over-fitting and tend to forget less throughout the training process\cite{tirumala2022memorization}.

% Machine learning turns out that machine learning models memorizes data (overfitting). Privacy attack 

% privacy attack > membership inference attack

% % data reconstruction attacks. 

% % private machine learning %(Jayaraman & Evans, 2019; Jagielski et al., 2020; Zanella-Béguelin et al., 2020; Nasr et al., 2021; Huang et al., 2022; Nasr et al., 2023; Steinke et al., 2023). 




\todo{Challenge I. Benchmark how to determine member and non-member. }


\todo{Challenge II. code syntax. }


challenge: verbatim instances


\textbf{Contribution.} We summarize our contributions as follows:





